\scanpage{197}
\mspage{6.}
correspondant à des $(P_i,G_i)$.

Dans le cas où l'objet final de $\mathscr R_S$ est connexe, on peut prendre les $P_i$ connexes, donc les morphismes de transition $G_j-G_i$ sont surjectifs, et on trouve $Z_i\simeq Z_j/H_{\unsure{ij|j}}$ \ldots. D'ailleurs, les $Z_i/G_i$ seront connexes \ldots. D'autre part, relativement à un tel système projectif, on peut répéter un énoncé comme prop. 4, soit

\uline{Prop} 5 \ldots..

Mais \deleted{\unsure{par les}} \uline{partant} de $\mathscr R$, et construisant un système projectif $(Z_i,G_i)$, \deleted{d'où} grâce à celui-ci un nouveau \deleted{champ} $(\mathscr R',r')$, et $\mathscr R'\hookrightarrow\mathscr R^{\deleted{\unsure{\prime|r|\bullet}}}$, il n'est pas vrai en général, \inserted{a priori,} que $\mathscr R'\hookrightarrow\mathscr R$ soit une équivalence : \uline{Le champ} $\mathscr R$ \uline{n'est pas nécessairement} ``\uline{engendré}'' par ses sections globales. Donc pour définir un $\mathscr R$ en \uline{général}, il faut se localiser et recoller \ldots. On va donc décrire un procédé de recollement qui nous permettra de construire les données de ramification dont nous aurons besoin.

Notions auxiliaires : Si $(\mathscr R,r)$, $(\mathscr R',r')$ sont deux données de ram. sur $S$, un hom. de l'une sur l'autre est un couple $(\varphi,\lambda)$, avec $\varphi:\mathscr R\longrightarrow\mathscr R'$ \deleted{\unsure{un}} foncteur cartésien, exact \inserted{sur les fibres,} et $\lambda$ un \uline{isomorphisme} de foncteurs cartésiens $r'\varphi\simeq\deleted{\unsure{\varphi}}r$. \deleted{Si $\mathscr R$ est défini} De cette façon,
% archive sheet number: 196
